\chapter{Cold Intolerance (Feeling Cold)}

\section*{Definition}
An exaggerated sensitivity to cold temperatures or a persistent sensation of feeling cold even in normal environmental conditions, often with visible vasoconstriction of the extremities.

\section*{Pathophysiology}
Cold intolerance results from altered thermoregulation, reduced metabolic heat production, or impaired vasomotor control. Oestrogen influences hypothalamic thermoregulatory set-point; declining oestrogen destabilises the thermoregulatory centre, causing both hot flushes and cold spells. Reduced basal metabolic rate (age-related and thyroid-related) decreases heat production. Vasomotor symptoms alternate between vasodilation (hot flushes) and vasoconstriction (cold flushes) as the hypothalamus attempts to regulate temperature.

\section*{Causes}
\begin{itemize}
  \item Perimenopausal vasomotor instability (cold flushes)
  \item Hypothyroidism (reduced metabolic rate)
  \item Iron deficiency anaemia (impaired oxygen delivery)
  \item Peripheral arterial disease
  \item Raynaud phenomenon or disease
  \item Low body weight or inadequate caloric intake
  \item Anorexia nervosa (loss of insulating fat)
  \item Vitamin B12 deficiency (impaired myelination)
  \item Diabetes mellitus (peripheral neuropathy)
  \item Medications (beta-blockers, ergot derivatives)
\end{itemize}

\section*{When to See a Doctor}
See your doctor if:
\begin{itemize}
  \item Severe or worsening symptoms
  \item Persistent cold intolerance with fatigue, weight gain, constipation, dry skin, or colour changes in fingers/toes (blue/white)
  \item Accompanied by other concerning signs
\end{itemize}

\section*{Related Symptoms}
\begin{itemize}
  \item Fatigue
  \item Cold hands and feet
  \item Weight gain
  \item Dry skin
  \item Hot flushes
\end{itemize}

\textbf{Important:} If this symptom persists, your doctor will check for serious underlying causes.

\section*{Self-Care / Home Management}
\begin{itemize}
  \item Rest and avoid activities that make it worse.
  \item Maintain adequate hydration and a balanced diet.
  \item Over-the-counter remedies may provide symptomatic relief (consult a pharmacist or doctor).
  \item Monitor symptoms and seek professional advice if they persist or worsen.
  \item Avoid self-medicating with prescription medications without medical guidance.
\end{itemize}

\section*{How Your Doctor Will Evaluate You}
\begin{itemize}
  \item A thorough history and physical examination are the first steps in evaluation.
  \item Laboratory tests (blood work, urinalysis) may be ordered based on clinical suspicion.
  \item Imaging studies (X-ray, ultrasound, CT, MRI) may be indicated when structural pathology is suspected.
  \item Specialist referral is arranged when the diagnosis remains uncertain or requires specific expertise.
\end{itemize}

\section*{Treatment Options}
\begin{itemize}
  \item Treatment targets the underlying cause identified through diagnostic evaluation.
  \item Supportive care and symptom management are provided while awaiting definitive therapy.
  \item Medications, lifestyle modifications, and physical therapies may be prescribed as appropriate.
  \item Follow-up ensures treatment effectiveness and allows timely adjustments to the care plan.
\end{itemize}

\clearpage